\subsection*{A.10 source-fixed interpolation closes the provisional A.11 angular-derivative gates (22 October 2026)}
\label{subsec:a11-A10-fixed-stage-closure}
The original Appendix A.10--A.11 construction fixes all stage boundaries
independently of \(\eta\). For fixed \(T_f=512\) and
\(0<\lambda\le1/100\), the A.11 first bump center follows
the interpolation endpoint by
\(L_1=30\log(1/\lambda)-3\), so \(\partial_\eta L_1=0\).
The moving-endpoint formula of the preceding section is still valid
for hypothetical moving stages, but its \(L_{1,\eta}\) term vanishes
\emph{for this source schedule}.

Write \(a=1-\lambda\), \(J_0=\log(1+\eta^2)\), and
take the source's exact smooth step \(\sigma\) from (A.5).
If the pre-interpolation moment ratio \(r_p=I/(XH)\) has its
source-factorized \(\eta\)-independence, the integrating factor in (A.10)
and \(\sigma(1-t)=1-\sigma(t)\) give
\[
r_f=\frac{2r_p}{1+\eta^2}e^{-aT_f}+
 \int_0^{T_f} e^{-av}
 \left(\frac2{1+\eta^2}\right)^{\sigma(v/T_f)}\,dv.
\]
Differentiation under this genuine finite integral gives
\[
\boxed{
r_{f,\eta}=-\frac{2\eta}{1+\eta^2}
\left[\frac{2r_p}{1+\eta^2}e^{-aT_f}+
\int_0^{T_f}e^{-av}\sigma(v/T_f)
\left(\frac2{1+\eta^2}\right)^{\sigma(v/T_f)}\,dv\right].
}
\]
For \(0\le r_p\le20/17\), \(a\ge99/100\), \(T_f=512\),
the exact step satisfies
\(\sigma(1/4)=1/(1+e^{128/9})<10^{-6}\);
the first 16 positive Taylor terms prove \(e^{128/9}>10^6\).
The first 15 terms give \(e^{30}>10^{10}\), so
\(e^{-aT_f/4},e^{-aT_f}<10^{-10}\).
Splitting the integral at \(T_f/4\) proves
\[
\boxed{|r_{f,\eta}|<
\frac{40}{17}\,10^{-10}
+\frac{200}{99}(10^{-6}+10^{-10})<3\cdot10^{-6}.}
\]
On the source-fixed \(-\lambda\) plateau the exact mismatch
\(\Delta_I=1/a-r(y_1)\) satisfies
\[
\Delta_I=-e^{-aL_1}(r_f-1/a),\qquad
\Delta_{I,\eta}=-e^{-aL_1}r_{f,\eta}.
\]
Assuming the already established \(0\le r_f\le20/17\)
(which follows from its incoming bound and the source slope barrier),
\(e^{-aL_1}\le e^3\lambda^{30a}\le21\lambda^{29}\) yields
\[
|\Delta_I|\le(21/95)\lambda^{28},\quad
|\Delta_{I,\eta}|<63\cdot10^{-6}\lambda^{29}.
\]
Together with the existing genuine A.11 small-branch inverse and
Jacobian gates \(\|c\|_\infty\le3|\Delta_I|\) and
\(\|c_\eta\|_\infty\le2|\Delta_{I,\eta}|\), this gives the
source-conditional sharpened bounds
\[
\boxed{\|c\|_\infty\le(63/95)\lambda^{28}<\lambda^{28},
\qquad\|c_\eta\|_\infty<\lambda^{28}/100000.}
\]
These finite source-stage deductions replace the provisional
\(|D_\eta r(a(\eta),\eta)|\le1/4\) and
\(|L_{1,\eta}|\le1/8\) assumptions of Edition III.
They recover and sharpen earlier A.10 sensitivity analyses
already contained in the cumulative manuscript, not a new
unconditional A.19 theorem.

The original source theorem has \emph{forcing}. The full ordered
parameter selection, same-profile A.19/pressure terminal gauges,
exterior stress cones, Appendix B contraction, unforced startup,
zero-force PDE residual, and complete Lean kernel proof remain
OPEN/HOLD. The finite rational program only audits the displayed
bounds and synthetic controls.
